\section{Nederlandse overheidscontext}
\subsection{Stelsel van Basisregistraties}
De historische eisen voor basisregistraties verbinden registratie met wettelijke verankering, verantwoordelijkheden, kwaliteit en terugmelding. De kwaliteitsbaseline werkt stelselbrede kwaliteitszorg verder uit. Deze bronnen zijn relevant voor MDM omdat zij laten zien dat gedeeld gebruik van kerngegevens niet losstaat van institutionele afspraken. De kwalificatie ``masterdata'' verleent op zichzelf echter geen authenticiteit of bevoegdheid. \parencite{S25,S26}\claim{RC45}

Het Stelsel kan niet zonder nadere analyse als één commerciële MDM-hub worden beschreven. Een vergelijking moet onderscheid maken tussen de verantwoordelijke registratie, de gegevensverstrekking en een eventuele samengestelde afnemersweergave. Dit is een toepassingsinterpretatie, geen claim dat alle basisregistraties op dezelfde technische wijze zijn ingericht. De paper onderzoekt geen operationele keten tussen registraties.

De BAG-catalogus biedt een concrete documentatiecontext met onder meer nummeraanduiding, verblijfsobject en pand. Een overeenkomstig label in een andere registratie bewijst geen identiteit. Voor een toekomstige vergelijking zijn brondefinities, modelrelaties en tijdsbetekenis nodig. Hier dient BAG uitsluitend om duidelijk te maken welke soorten documenten en begrippen bij een overheidsgerichte MDM-analyse moeten worden onderscheiden. Er worden geen niet-gecontroleerde BRP-, BRK- of WOZ-relaties ingevuld. \parencite{S49}\claim{RC46}

\subsection{Nederlandse modellen, begrippen en metadata}
MIM ordent conceptuele en logische informatiemodellering; NL-SBB ordent begripsbeschrijvingen; TOOI biedt een kennismodel voor overheidsinformatie. NEN 3610 vormt een afzonderlijke normatieve context voor geo-informatie. DCAT-AP-NL adresseert catalogusmetadata. Dit zijn verschillende abstractieniveaus en toepassingsgebieden, zodat een gezamenlijk gebruik niet mag worden teruggebracht tot een keuze voor één universele modeltaal. \parencite{S13,S14,S15,S17,S16}\claim{RC47}

Voor MDTO is het relevante vraagstuk duurzame toegankelijkheid en metagegevens van informatieobjecten. Het kan daardoor raken aan archivering en historie, maar is niet zonder meer een vervanging voor een masterdata- of begrippenmodel. BOMOS gaat juist over beheer van standaarden. De twee documenten delen een beheerthema zonder noodzakelijk hetzelfde beheerobject te hebben. \parencite{S24,S23}

\subsection{Versie, lijststatus en toepasselijkheid}
Bij de officiële lijstcontrole op 9 september 2026 is MIM 1.2 aanbevolen. De verplichte lijst vermeldt NL-SBB 1.0, OpenAPI Specification 3.0, REST API Design Rules 2.2 en NL GOV CloudEvents 1.1. DCAT-AP-NL 3.0 staat bij de standaarden in behandeling. Ook TOOI en MDTO staan op de geraadpleegde lijst in behandeling. Dit zijn waarnemingen van de officiële registraties op de peildatum, geen zelfstandig oordeel over iedere organisatie of gegevensuitwisseling. \parencite{F01,F02,F03}\claim{RC48}

De technische analyse gebruikt onder meer OpenAPI 3.1.1. Dat is een specificatieversie en niet automatisch de versie die op de Nederlandse lijst is vastgesteld. Formele status en technische documentversie worden daarom in de matrix gescheiden. De paper leidt uit de aanwezigheid op een lijst geen onvoorwaardelijke wettelijke verplichting voor alle MDM-functionaliteit af. \parencite{S19,F01,F03}

\subsection{Federatief Datastelsel en Europese kaders}
Het geraadpleegde FDS-afwegingskader vormt een ingang naar voorwaarden voor toelating van aanbod. De beschikbare hoofdpagina is echter onvoldoende om alle operationele afspraken van FDS inhoudelijk te reconstrueren. De studie positioneert FDS daarom als relevante Nederlandse onderzoekscontext en laat gedetailleerde vergelijking met federatieve MDM-patronen open. \parencite{S27}\claim{RC49}

EIF ordent interoperabiliteit over meerdere niveaus. De Interoperable Europe Act vormt een juridisch kader voor interoperabiliteit in de Europese publieke sector. SEMIC Core Vocabularies bieden herbruikbare kernmodellen, bijvoorbeeld voor personen, organisaties en locaties. Hun rol verschilt: een juridisch of bestuurlijk kader is geen datamodel, en een kernvocabulaire beslist niet automatisch over de volledige betekenis van een nationale registratierelatie. \parencite{S43,S42,S44}\claim{RC50}

FAIR benadrukt vindbaarheid, toegankelijkheid, interoperabiliteit en herbruikbaarheid, inclusief aandacht voor machinegericht gebruik. De principes ondersteunen daarmee het belang van metadata en expliciete verwijzingen. Zij verplichten niet alle gegevens openbaar te maken en bewijzen geen bepaalde MDM-architectuur. \parencite{L11}\claim{RC51}

\section{Masterdata en Data Products}
Dataproductdenken legt nadruk op een terugkerende informatiebehoefte, beheer en consumptie. Een dataproduct kan masterdata aanbieden, maar ook gegevens uit andere categorieën. De relatie is daarom geen eenvoudige vervanging van de ene term door de andere. Productdenken betreft het georganiseerde aanbod; masterdata betreft een gegevens- en beheervraagstuk dat daarbinnen kan liggen. \parencite{L12}\claim{RC52}

De gedachte dat masterdata als product kunnen worden behandeld is ouder dan de recente Data Mesh-discussie. Otto en Ofner noemen dit expliciet en verbinden er productbeschrijvingen en zicht op bronnen en afnemers aan. Een nieuw label is dus onvoldoende om een nieuwe capability of wetenschappelijke bijdrage te claimen. Wel blijft onderzoek nodig naar de praktische invulling en verschillen tussen oudere en recente productbenaderingen. \parencite{L20}\claim{RC53}

DPROD biedt een kandidaatvocabulaire voor productmetadata, waaronder eigenaar, doel en relaties met datasets en diensten. De hier geraadpleegde versie is 1.0 Beta 1. De bron ondersteunt het bestaan van een productspecificatie, niet dat deze versie iedere MDM-behoefte afdekt of dat een specifiek product ermee conformeert. Samenhang met DCAT-profielen en metadataregistries vraagt een vergelijking van resources en verantwoordelijkheden. \parencite{S51,S47,S08}\claim{RC54}

\section{Masterdata en federatieve data-architecturen}
De grijze-literatuurreview van Goedegebuure et al. synthetiseert praktijkmateriaal over Data Mesh. De beschreven principes betreffen domeinverantwoordelijkheid, data als product, self-service en federatieve governance. De geraadpleegde auteursversie is geen door deze paper uitgevoerd praktijkonderzoek; de onderliggende professionele claims worden niet automatisch bewijs van gerealiseerde voordelen. \parencite{L23}\claim{RC55}

Federatie betreft bovendien meerdere dimensies: verdeling van gegevens, verdeling van besluitrechten en de afspraken waarmee partijen samenwerken. Een centraal metadataoverzicht kan met decentraal gegevensbeheer samengaan. Daarom is een tegenstelling tussen ``centraal MDM'' en ``decentraal Data Mesh'' analytisch te grof. De patroonbeschrijvingen en recente architectuurvergelijkingen bevatten verschillende combinaties; de optimale keuze volgt daar niet algemeen uit. \parencite{L20,L23,L24}\claim{RC56}

Gieß en Hutterer ordenen data platforms, spaces, meshes en fabrics onder meer naar governance en toepassingsgebied. Hun synthese is bruikbaar als vergelijkingskader, maar een classificatie op dat niveau bepaalt geen formele consistentiegaranties of lokale beheerlast. Beheerautonomie kan extra afstemming vragen; centralisatie kan bevoegdheden bundelen. Beide zijn afwegingsdimensies die in een concrete context moeten worden onderzocht. \parencite{L24}

\section{API-first, events en data contracts}
OpenAPI beschrijft interfaces, de Nederlandse REST-regels stellen ontwerpafspraken en Digikoppeling/FSC adresseert onderdelen van gegevensuitwisseling en verbindingen. Het NL GOV-profiel voor CloudEvents beschrijft gebeurtenissen binnen een uitwisselingscontext. Deze instrumenten kunnen distributie en verandering ondersteunen, maar definiëren niet vanzelf de domeinbetekenis van een merge, correctie of intrekking. \parencite{S19,S18,S20,S21,S22}\claim{RC57}

``API-first'' en ``event-driven'' worden hier als ontwerporiëntaties onderzocht, niet als bewezen opvolgers van MDM. Hetzelfde geldt voor data contracts. De term kan een technisch schema, een kwaliteitsafspraak of een breder productcontract aanduiden. DPROD, catalogusprofielen en interfacespecificaties bieden aangrenzende constructies, maar het huidige corpus rechtvaardigt geen universele definitie of rangorde van data-contractbenaderingen. Dat is een beperking van deze synthese. \parencite{S51,S47,S19}\claim{RC58}

Een event dat een wijziging meldt, hoeft niet de volledige nieuwe toestand te bevatten. Voor interpretatie moet daarom worden onderzocht welke identifiers, versiegegevens en verwijzingen het gekozen profiel en de concrete toepassing leveren. Dit is een analysevraag op het grensvlak van eventing en historie; er worden geen niet-uitgevoerde prestatietests van eventverwerking gerapporteerd. \parencite{S22,S37}

\section{Masterdata en AI}
AI kan zowel hulpmiddel bij gegevensbeheer als afnemer van gegevens zijn. Entity-resolutiononderzoek behandelt algoritmische matching en de xLP-publicatie illustreert uitlegbaarheid rond voorspelde relaties in een MDM-context. Het bestaan van deze toepassingen zegt niets over een universele vervanging van governance of over de betrouwbaarheid van ieder afgeleid verband. De geraadpleegde xLP-tekst is een in 2024 gedeponeerde auteursversie met andere jaaraanduidingen in het document; deze paper gebruikt haar als voorbeeld, niet om een chronologische prioriteitsclaim te maken. \parencite{L29,L31}\claim{RC59}

Retrieval-Augmented Generation combineert in Lewis et al. generatie met opgehaalde informatie. De studie rapporteert resultaten voor haar eigen taken en modellen. Zij bewijst niet dat een MDM-product met ontologieën automatisch betere antwoorden produceert. Knowledge graphs en expliciete metadata kunnen informatie structureren en ontsluiten, maar onjuiste bronnen, verkeerde mappings of een verkeerde selectiecontext blijven mogelijke fouten. \parencite{L16,L13,S33}\claim{RC60}

Semantic grounding wordt hier opgevat als het verbinden van een interpretatie aan controleerbare concepten en bronnen. Dat is een analytische werkbetekenis, geen claim dat het corpus een algemeen geaccepteerde definitie levert. Het onderscheid tussen normatieve definitie, gezaghebbende registratie, overige registratie en afleiding kan bij AI-consumptie relevant zijn. Hoe die categorieën machineleesbaar worden weergegeven en of zij fouten verminderen, blijft een open onderzoeksvraag. Provenance is daarvoor een mogelijke bouwsteen, geen bewijs van effectiviteit. \parencite{S33,S25,L16}\claim{RC61}

Voor AI-agents als zelfstandige ketenafnemers is de specifieke empirische dekking in deze verzameling beperkt. Claims over veilige autonome besluitvorming, hallucinatie-eliminatie of gegarandeerde interpretatie worden daarom niet gedaan. De noodzakelijke vervolgstudie moet toegang, bronversie, modelgedrag, taakcontext en foutkosten samen behandelen. Deze onderzoeksagenda volgt uit de onderscheiden bewijsgrenzen, niet uit een uitgevoerd experiment.
